\documentclass[10pt,a4paper]{amsart}
\usepackage{amsmath,amssymb,amsthm,graphicx}
% Retained preamble for the Stacks Project chapter snapshot in this directory.
% The chapter's own preamble.tex is a copy of the shared Stacks Project preamble; the block
% below is sliced MECHANICALLY, line for line, from it. Nothing was added, removed or
% redefined; the wrapper class, the xr-hyper/hyperref wiring and the externaldocument list
% are omitted because the assembler supplies the document class.
%
% lines taken from preamble.tex: 3, 5, 9, 11, 12, 156, 157, 158, ... (51 lines total)
%
% package coverage restored from the chapter's own preamble (2026-10-13): xy, xy-2cell

\usepackage[all]{xy}\usepackage{verbatim}
\xyoption{2cell}
\UseAllTwocells

\newenvironment{reference}{\comment}{\endcomment}
\newenvironment{slogan}{\comment}{\endcomment}
\newenvironment{history}{\comment}{\endcomment}
\theoremstyle{plain}
\newtheorem{theorem}[subsection]{Theorem}
\newtheorem{proposition}[subsection]{Proposition}
\newtheorem{lemma}[subsection]{Lemma}
\theoremstyle{definition}
\newtheorem{definition}[subsection]{Definition}
\newtheorem{example}[subsection]{Example}
\newtheorem{exercise}[subsection]{Exercise}
\newtheorem{situation}[subsection]{Situation}
\theoremstyle{remark}
\newtheorem{remark}[subsection]{Remark}
\newtheorem{remarks}[subsection]{Remarks}
\numberwithin{equation}{subsection}
\def\lim{\mathop{\mathrm{lim}}\nolimits}
\def\colim{\mathop{\mathrm{colim}}\nolimits}
\def\Spec{\mathop{\mathrm{Spec}}}
\def\Hom{\mathop{\mathrm{Hom}}\nolimits}
\def\Ext{\mathop{\mathrm{Ext}}\nolimits}
\def\SheafHom{\mathop{\mathcal{H}\!\mathit{om}}\nolimits}
\def\SheafExt{\mathop{\mathcal{E}\!\mathit{xt}}\nolimits}
\def\Sch{\mathit{Sch}}
\def\Mor{\mathop{\mathrm{Mor}}\nolimits}
\def\Ob{\mathop{\mathrm{Ob}}\nolimits}
\def\Sh{\mathop{\mathit{Sh}}\nolimits}
\def\NL{\mathop{N\!L}\nolimits}
\def\CH{\mathop{\mathrm{CH}}\nolimits}
\def\proetale{{pro\text{-}\acute{e}tale}}
\def\etale{{\acute{e}tale}}
\def\QCoh{\mathit{QCoh}}
\def\Ker{\mathop{\mathrm{Ker}}}
\def\Im{\mathop{\mathrm{Im}}}
\def\Coker{\mathop{\mathrm{Coker}}}
\def\Coim{\mathop{\mathrm{Coim}}}
\DeclareMathSymbol{\boxtimes}{\mathbin}{AMSa}{"02}
\def\QCohstack{\mathcal{QC}\!\mathit{oh}}
\def\Cohstack{\mathcal{C}\!\mathit{oh}}
\def\Spacesstack{\mathcal{S}\!\mathit{paces}}
\def\Quotfunctor{\mathrm{Quot}}
\def\Hilbfunctor{\mathrm{Hilb}}
\def\Curvesstack{\mathcal{C}\!\mathit{urves}}
\def\Polarizedstack{\mathcal{P}\!\mathit{olarized}}
\def\Complexesstack{\mathcal{C}\!\mathit{omplexes}}
\def\Pic{\mathop{\mathrm{Pic}}\nolimits}
\def\Picardstack{\mathcal{P}\!\mathit{ic}}
\def\Picardfunctor{\mathrm{Pic}}
\def\Deformationcategory{\mathcal{D}\!\mathit{ef}}
\usepackage{cleveref}
\crefname{originallemma}{original lemma}{original lemmas}
\Crefname{originallemma}{Original lemma}{Original lemmas}
\makeatletter
\newlabel{algebra-lemma-universally-bijective}{{09EF}{source}{}{sourcepointer-1}{}}
\newlabel{algebra-lemma-universally-bijective@cref}{{[originallemma][0][]09EF}{[1][0][]source}{}{}{}}
\makeatother
\crefname{originallemma}{original lemma}{original lemmas}
\Crefname{originallemma}{Original lemma}{Original lemmas}
\makeatletter
\newlabel{groupoids-lemma-integral-over-invariants}{{03BJ}{source}{}{sourcepointer-2}{}}
\newlabel{groupoids-lemma-integral-over-invariants@cref}{{[originallemma][0][]03BJ}{[1][0][]source}{}{}{}}
\makeatother
\crefname{originallemma}{original lemma}{original lemmas}
\Crefname{originallemma}{Original lemma}{Original lemmas}
\makeatletter
\newlabel{groupoids-lemma-invariants-base-change}{{03BK}{source}{}{sourcepointer-3}{}}
\newlabel{groupoids-lemma-invariants-base-change@cref}{{[originallemma][0][]03BK}{[1][0][]source}{}{}{}}
\makeatother
\crefname{originallemma}{original lemma}{original lemmas}
\Crefname{originallemma}{Original lemma}{Original lemmas}
\makeatletter
\newlabel{groupoids-lemma-points}{{03BL}{source}{}{sourcepointer-4}{}}
\newlabel{groupoids-lemma-points@cref}{{[originallemma][0][]03BL}{[1][0][]source}{}{}{}}
\makeatother
\crefname{originallemma}{original lemma}{original lemmas}
\Crefname{originallemma}{Original lemma}{Original lemmas}
\makeatletter
\newlabel{groupoids-quotients-lemma-base-change-quotient-stack}{{0DTG}{source}{}{sourcepointer-5}{}}
\newlabel{groupoids-quotients-lemma-base-change-quotient-stack@cref}{{[originallemma][0][]0DTG}{[1][0][]source}{}{}{}}
\makeatother
\crefname{originallemma}{original lemma}{original lemmas}
\Crefname{originallemma}{Original lemma}{Original lemmas}
\makeatletter
\newlabel{groupoids-section-finite-flat}{{03BE}{source}{}{sourcepointer-6}{}}
\newlabel{groupoids-section-finite-flat@cref}{{[originallemma][0][]03BE}{[1][0][]source}{}{}{}}
\makeatother
\crefname{originallemma}{original lemma}{original lemmas}
\Crefname{originallemma}{Original lemma}{Original lemmas}
\makeatletter
\newlabel{lemma-categorical-quotient-compare}{{lemma-categorical-quotient-compare}{source}{}{sourcepointer-7}{}}
\newlabel{lemma-categorical-quotient-compare@cref}{{[originallemma][0][]lemma-categorical-quotient-compare}{[1][0][]source}{}{}{}}
\makeatother
\crefname{originallemma}{original lemma}{original lemmas}
\Crefname{originallemma}{Original lemma}{Original lemmas}
\makeatletter
\newlabel{lemma-quotient-compare}{{lemma-quotient-compare}{source}{}{sourcepointer-8}{}}
\newlabel{lemma-quotient-compare@cref}{{[originallemma][0][]lemma-quotient-compare}{[1][0][]source}{}{}{}}
\makeatother
\crefname{originallemma}{original lemma}{original lemmas}
\Crefname{originallemma}{Original lemma}{Original lemmas}
\makeatletter
\newlabel{lemma-well-nigh-affine}{{lemma-well-nigh-affine}{source}{}{sourcepointer-9}{}}
\newlabel{lemma-well-nigh-affine@cref}{{[originallemma][0][]lemma-well-nigh-affine}{[1][0][]source}{}{}{}}
\makeatother
\crefname{originallemma}{original lemma}{original lemmas}
\Crefname{originallemma}{Original lemma}{Original lemmas}
\makeatletter
\newlabel{morphisms-lemma-integral-universally-closed}{{01WM}{source}{}{sourcepointer-10}{}}
\newlabel{morphisms-lemma-integral-universally-closed@cref}{{[originallemma][0][]01WM}{[1][0][]source}{}{}{}}
\makeatother
\crefname{originallemma}{original lemma}{original lemmas}
\Crefname{originallemma}{Original lemma}{Original lemmas}
\makeatletter
\newlabel{schemes-lemma-morphism-into-affine}{{01I1}{source}{}{sourcepointer-11}{}}
\newlabel{schemes-lemma-morphism-into-affine@cref}{{[originallemma][0][]01I1}{[1][0][]source}{}{}{}}
\makeatother
\crefname{originallemma}{original lemma}{original lemmas}
\Crefname{originallemma}{Original lemma}{Original lemmas}
\makeatletter
\newlabel{stacks-morphisms-lemma-compose-after-separated}{{050M}{source}{}{sourcepointer-12}{}}
\newlabel{stacks-morphisms-lemma-compose-after-separated@cref}{{[originallemma][0][]050M}{[1][0][]source}{}{}{}}
\makeatother
\crefname{originallemma}{original lemma}{original lemmas}
\Crefname{originallemma}{Original lemma}{Original lemmas}
\makeatletter
\newlabel{stacks-morphisms-lemma-image-proper-is-proper}{{0CQK}{source}{}{sourcepointer-13}{}}
\newlabel{stacks-morphisms-lemma-image-proper-is-proper@cref}{{[originallemma][0][]0CQK}{[1][0][]source}{}{}{}}
\makeatother
\crefname{originallemma}{original lemma}{original lemmas}
\Crefname{originallemma}{Original lemma}{Original lemmas}
\makeatletter
\newlabel{stacks-morphisms-lemma-points-presentation}{{0DTZ}{source}{}{sourcepointer-14}{}}
\newlabel{stacks-morphisms-lemma-points-presentation@cref}{{[originallemma][0][]0DTZ}{[1][0][]source}{}{}{}}
\makeatother
\crefname{originallemma}{original lemma}{original lemmas}
\Crefname{originallemma}{Original lemma}{Original lemmas}
\makeatletter
\newlabel{stacks-morphisms-lemma-surjection-from-quasi-compact}{{050X}{source}{}{sourcepointer-15}{}}
\newlabel{stacks-morphisms-lemma-surjection-from-quasi-compact@cref}{{[originallemma][0][]050X}{[1][0][]source}{}{}{}}
\makeatother
\crefname{originallemma}{original lemma}{original lemmas}
\Crefname{originallemma}{Original lemma}{Original lemmas}
\makeatletter
\newlabel{stacks-morphisms-lemma-universally-injective-local}{{0DTN}{source}{}{sourcepointer-16}{}}
\newlabel{stacks-morphisms-lemma-universally-injective-local@cref}{{[originallemma][0][]0DTN}{[1][0][]source}{}{}{}}
\makeatother
\title{The well-nigh-affine invariant-ring moduli space}
\author{The Stacks Project Authors}
\date{stacks-project:stacks-more-morphisms.tex:local-snapshot-20261012; GNU FDL 1.2 or later}
\begin{document}
\maketitle
\noindent\textit{Editorial note.} Statement span is lines 3424-3433 of sources/stacks-stacks-more-morphisms/pinned-original-source.tex; proof span is lines 3435-3514, the complete author proof (78 body lines).\par\medskip
\noindent\textit{Editorial note.} Navigation stubs: 16 cited-but-not-included labels are supplied as non-load-bearing pointers with their pinned tags where the snapshot tag index has them.\par\medskip
\noindent\textit{Editorial note.} extra\_packages is empty on purpose: the retained preamble already loads xy with options and a second declaration clashes in the assembler wrapper.\par\medskip
\noindent\textit{Editorial note.} Text is retained verbatim, including the author's own typos and placeholders.\par\medskip
\noindent\textit{Editorial note.} Queue index 12 of queue\_stacks-stacks-more-morphisms.json (pb=3365, stmt 3424-3433, proof 3435-3514).\par\medskip
\noindent\textit{Editorial note.} stacks-stacks-more-morphisms wave smm01 (2026-10-14): fresh adjudication accept per lane-s3/smm\_judgments.md; pool per Chief 10:5x open-pool ruling (net 59 with-proof rows, bloodline lane-s3 by anchor audit, v2 true-key namespace stacks:stacks-more-morphisms.tex, row lemma-deform-property-fp-over-ft avoided on cross-chapter key collision); 60 with-proof rows judged 45 accept / 14 reject / 1 avoid; patch-sd-envs contract spans scanned.\par\medskip
\noindent\textit{Allowed standard prerequisites (editorial).}\begin{itemize}
\item Explicit prerequisite boundary. The cited inputs are prerequisites; the argument in the retained span is not replaced by any cited result.
\end{itemize}\medskip
\ifdefined\hypertarget\hypertarget{sourcepointer-1}{}\fi
\noindent\textit{Source reference.} Original lemma 09EF: sources/stacks-stacks-more-morphisms/pinned-original-source.tex, source lines 3514--3514. Cited by the retained statement or proof but not part of the collected spans; supplied only as a navigation pointer. The argument does not depend on its content.\par\medskip
\ifdefined\hypertarget\hypertarget{sourcepointer-2}{}\fi
\noindent\textit{Source reference.} Original lemma 03BJ: sources/stacks-stacks-more-morphisms/pinned-original-source.tex, source lines 3514--3514. Cited by the retained statement or proof but not part of the collected spans; supplied only as a navigation pointer. The argument does not depend on its content.\par\medskip
\ifdefined\hypertarget\hypertarget{sourcepointer-3}{}\fi
\noindent\textit{Source reference.} Original lemma 03BK: sources/stacks-stacks-more-morphisms/pinned-original-source.tex, source lines 3514--3514. Cited by the retained statement or proof but not part of the collected spans; supplied only as a navigation pointer. The argument does not depend on its content.\par\medskip
\ifdefined\hypertarget\hypertarget{sourcepointer-4}{}\fi
\noindent\textit{Source reference.} Original lemma 03BL: sources/stacks-stacks-more-morphisms/pinned-original-source.tex, source lines 3514--3514. Cited by the retained statement or proof but not part of the collected spans; supplied only as a navigation pointer. The argument does not depend on its content.\par\medskip
\ifdefined\hypertarget\hypertarget{sourcepointer-5}{}\fi
\noindent\textit{Source reference.} Original lemma 0DTG: sources/stacks-stacks-more-morphisms/pinned-original-source.tex, source lines 3514--3514. Cited by the retained statement or proof but not part of the collected spans; supplied only as a navigation pointer. The argument does not depend on its content.\par\medskip
\ifdefined\hypertarget\hypertarget{sourcepointer-6}{}\fi
\noindent\textit{Source reference.} Original lemma 03BE: sources/stacks-stacks-more-morphisms/pinned-original-source.tex, source lines 3514--3514. Cited by the retained statement or proof but not part of the collected spans; supplied only as a navigation pointer. The argument does not depend on its content.\par\medskip
\ifdefined\hypertarget\hypertarget{sourcepointer-7}{}\fi
\noindent\textit{Source reference.} Original lemma lemma-categorical-quotient-compare: sources/stacks-stacks-more-morphisms/pinned-original-source.tex, source lines 3247--3252. Cited by the retained statement or proof but not part of the collected spans; supplied only as a navigation pointer. The argument does not depend on its content.\par\medskip
\ifdefined\hypertarget\hypertarget{sourcepointer-8}{}\fi
\noindent\textit{Source reference.} Original lemma lemma-quotient-compare: sources/stacks-stacks-more-morphisms/pinned-original-source.tex, source lines 3217--3224. Cited by the retained statement or proof but not part of the collected spans; supplied only as a navigation pointer. The argument does not depend on its content.\par\medskip
\ifdefined\hypertarget\hypertarget{sourcepointer-9}{}\fi
\noindent\textit{Source reference.} Original lemma lemma-well-nigh-affine: sources/stacks-stacks-more-morphisms/pinned-original-source.tex, source lines 3339--3346. Cited by the retained statement or proof but not part of the collected spans; supplied only as a navigation pointer. The argument does not depend on its content.\par\medskip
\ifdefined\hypertarget\hypertarget{sourcepointer-10}{}\fi
\noindent\textit{Source reference.} Original lemma 01WM: sources/stacks-stacks-more-morphisms/pinned-original-source.tex, source lines 3514--3514. Cited by the retained statement or proof but not part of the collected spans; supplied only as a navigation pointer. The argument does not depend on its content.\par\medskip
\ifdefined\hypertarget\hypertarget{sourcepointer-11}{}\fi
\noindent\textit{Source reference.} Original lemma 01I1: sources/stacks-stacks-more-morphisms/pinned-original-source.tex, source lines 3514--3514. Cited by the retained statement or proof but not part of the collected spans; supplied only as a navigation pointer. The argument does not depend on its content.\par\medskip
\ifdefined\hypertarget\hypertarget{sourcepointer-12}{}\fi
\noindent\textit{Source reference.} Original lemma 050M: sources/stacks-stacks-more-morphisms/pinned-original-source.tex, source lines 3514--3514. Cited by the retained statement or proof but not part of the collected spans; supplied only as a navigation pointer. The argument does not depend on its content.\par\medskip
\ifdefined\hypertarget\hypertarget{sourcepointer-13}{}\fi
\noindent\textit{Source reference.} Original lemma 0CQK: sources/stacks-stacks-more-morphisms/pinned-original-source.tex, source lines 3514--3514. Cited by the retained statement or proof but not part of the collected spans; supplied only as a navigation pointer. The argument does not depend on its content.\par\medskip
\ifdefined\hypertarget\hypertarget{sourcepointer-14}{}\fi
\noindent\textit{Source reference.} Original lemma 0DTZ: sources/stacks-stacks-more-morphisms/pinned-original-source.tex, source lines 3514--3514. Cited by the retained statement or proof but not part of the collected spans; supplied only as a navigation pointer. The argument does not depend on its content.\par\medskip
\ifdefined\hypertarget\hypertarget{sourcepointer-15}{}\fi
\noindent\textit{Source reference.} Original lemma 050X: sources/stacks-stacks-more-morphisms/pinned-original-source.tex, source lines 3514--3514. Cited by the retained statement or proof but not part of the collected spans; supplied only as a navigation pointer. The argument does not depend on its content.\par\medskip
\ifdefined\hypertarget\hypertarget{sourcepointer-16}{}\fi
\noindent\textit{Source reference.} Original lemma 0DTN: sources/stacks-stacks-more-morphisms/pinned-original-source.tex, source lines 3514--3514. Cited by the retained statement or proof but not part of the collected spans; supplied only as a navigation pointer. The argument does not depend on its content.\par\medskip
\begin{lemma}
\label{lemma-well-nigh-affine-moduli-space}
Let the algebraic stack $\mathcal{X}$ be well-nigh affine. There exists
a uniform categorical moduli space
$$
f : \mathcal{X} \longrightarrow M
$$
in the category of affine schemes. Moreover
$f$ is separated, quasi-compact, and a universal homeomorphism.
\end{lemma}

\begin{proof}
Write $\mathcal{X} = [U/R]$ with $(U, R, s, t, c)$ as in
Lemma \ref{lemma-well-nigh-affine}. Let $C$ be the ring of
$R$-invariant functions on $U$, see
Groupoids, Section \ref{groupoids-section-finite-flat}.
We set $M = \Spec(C)$. The $R$-invariant morphism
$U \to M$ corresponds to a morphism $f : \mathcal{X} \to M$ by
Lemma \ref{lemma-quotient-compare}.
The characterization of morphisms into affine schemes given in
Schemes, Lemma \ref{schemes-lemma-morphism-into-affine}
immediately guarantees that $\phi : U \to M$ is a categorical
quotient in the category of affine schemes. Hence $f$ is a
categorical moduli space in the category of affine schemes
(Lemma \ref{lemma-categorical-quotient-compare}).

\medskip\noindent
Since $\mathcal{X}$ is separated by Lemma \ref{lemma-well-nigh-affine}
we find that $f$ is separated by Morphisms of Stacks, Lemma
\ref{stacks-morphisms-lemma-compose-after-separated}.

\medskip\noindent
Since $U \to \mathcal{X}$ is surjective and since $U \to M$ is quasi-compact,
we see that $f$ is quasi-compact by Morphisms of Stacks, Lemma
\ref{stacks-morphisms-lemma-surjection-from-quasi-compact}.

\medskip\noindent
By Groupoids, Lemma \ref{groupoids-lemma-integral-over-invariants}
the composition
$$
U \to \mathcal{X} \to M
$$
is an integral morphism of affine schemes. In particular, it is
universally closed
(Morphisms, Lemma \ref{morphisms-lemma-integral-universally-closed}).
Since $U \to \mathcal{X}$ is surjective, it follows that $\mathcal{X} \to M$
is universally closed (Morphisms of Stacks, Lemma
\ref{stacks-morphisms-lemma-image-proper-is-proper}).
To conclude that $\mathcal{X} \to M$ is a universal homeomorphism,
it is enough to show that it is universally bijective, i.e.,
surjective and universally injective.

\medskip\noindent
We have $|\mathcal{X}| = |U|/|R|$ by
Morphisms of Stacks, Lemma \ref{stacks-morphisms-lemma-points-presentation}.
Thus $|f|$ is surjective and even bijective
by Groupoids, Lemma \ref{groupoids-lemma-points}.

\medskip\noindent
Let $C \to C'$ be a ring map. Let $(U', R', s', t', c')$ be
the base change of $(U, R, s, t, c)$ by $M' = \Spec(C') \to M$.
Setting $\mathcal{X}' = [U'/R']$, we observe that
$M' \times_M \mathcal{X} = \mathcal{X}'$ by
Quotients of Groupoids, Lemma
\ref{groupoids-quotients-lemma-base-change-quotient-stack}.
Let $C^1$ be the ring of $R'$-invariant functions on $U'$.
Set $M^1 = \Spec(C^1)$ and consider the diagram
$$
\xymatrix{
\mathcal{X}' \ar[d]^{f'} \ar[r] & \mathcal{X} \ar[dd]^f \\
M^1 \ar[d] \\
M' \ar[r] & M
}
$$
By Groupoids, Lemma \ref{groupoids-lemma-invariants-base-change} and
Algebra, Lemma \ref{algebra-lemma-universally-bijective}
the morphism $M^1 \to M'$ is a homeomorphism.
On the other hand, the previous paragraph applied to
$(U', R', s', t', c')$ shows that $|f'|$ is bijective.
We conclude that $f$ induces a bijection on points after any
base change by an affine scheme. Thus $f$ is universally injective
by Morphisms of Stacks, Lemma
\ref{stacks-morphisms-lemma-universally-injective-local}.

\medskip\noindent
Finally, we still have to show that $f$ is a uniform moduli space
in the category of affine schemes. This follows from the discussion
above and the fact that if the
ring map $C \to C'$ is flat, then $C' \to C^1$ is an isomorphism
by Groupoids, Lemma \ref{groupoids-lemma-invariants-base-change}.
\end{proof}

\end{document}
